\subsection{李群芽}

定义 5. K 上的李群芽是满足下列条件的系统 (G, e, θ, m)：

\begin{enumerate}
\def\labelenumi{(\roman{enumi})}
\item
  G 是 K 上的解析流形；
\item
  \(e \in G;\)
\item
  \(\theta\) 是从 G 到 G 的解析映射；
\item
  \(m\) 是从开子集 \(\Omega\)（\(\mathbf{G} \times \mathbf{G}\) 的开子集）到 \(\mathbf{G}\) 的解析映射；
\item
  对所有 \(g \in G\)，都有 \((e, g) \in \Omega\)、\((g, e) \in \Omega\) 以及 \(m(e, g) = m(g, e) = g\)；
\end{enumerate}

\[
g \in G, (g, \theta (g)) \in \Omega , (\theta (g), g) \in \Omega , m (g, \theta (g)) = m (\theta (g), g) = e;
\]

\begin{enumerate}
\def\labelenumi{(\roman{enumi})}
\setcounter{enumi}{6}
\tightlist
\item
  若 \(g, h, k\) 是 \(G\) 中满足 \((g, h) \in \Omega\)、\((h, k) \in \Omega\)、\((m(g, h), k) \in \Omega\) 和 \((g, m(h, k)) \in \Omega\) 的元素，则 \(m(m(g, h), k) = m(g, m(h, k))\)。
\end{enumerate}

e 称为李群芽的单位元。我们通常用 gh 代替 \(m(g, h)\)，并且（滥用记号）用 \(g^{-1}\) 代替 \(\theta(g)\)。

李群 G 是取 \(e, \theta, m\) 为显然选择的李群芽。设 G 是李群芽。则 \(ee^{-1} = e\)，即

\[
e ^ {- 1} = e.\tag{8}
\]

对所有 \(g\in G\)

\[
g = e g = ((g ^ {- 1}) ^ {- 1} g ^ {- 1}) g = (g ^ {- 1}) ^ {- 1} (g ^ {- 1} g) = (g ^ {- 1}) ^ {- 1} e,
\]

即

\[
(g ^ {- 1}) ^ {- 1} = g.\tag{9}
\]

G 的一个子集若在映射 \(g \mapsto g^{-1}\) 下不变，就称为对称的。流形 G 取点 e、映射 \(g \mapsto g^{-1}\) 以及映射 \((g, h) \mapsto hg\)，构成称为 G 的对立李群芽的李群芽 G \(^{\nu}\)。

若对于所有 \((g,h)\in \mathbf{G}\times \mathbf{G}\) 中满足 \(gh\) 有定义的元素，\(hg\) 也有定义且等于 \(gh\)，则称李群芽 G 为交换的。

设 \(G\) 是李群芽。由 \((g, h) \in G \times G\) 且 \(gh\) 有定义的元素所成的集合是 \((e, e)\) 的邻域。另一方面，映射 \((g, h) \mapsto gh\) 和 \(g \mapsto g^{-1}\) 是连续的。因此 \((gh)k = g(hk)\)，其中 \(g, h, k\) 充分接近 \(e\)。同样，\((h^{-1}g^{-1})(gh) = h^{-1}(eh) = h^{-1}h = e\)，其中 \(g, h\) 充分接近 \(e\)，于是右乘 \((gh)^{-1}\)，

\begin{enumerate}
\def\labelenumi{(\arabic{enumi})}
\setcounter{enumi}{9}
\tightlist
\item
\end{enumerate}

\((gh)^{-1} = h^{-1}g^{-1}\)，其中 \(g,h\) 充分接近 \(e\)。

命题 20. 设 G 是李群芽且 \(g \in G\)。存在 e 的开邻域 U 和 g 的开邻域 V，具有下列性质：

\begin{enumerate}
\def\labelenumi{(\alph{enumi})}
\item
  对所有 \(u \in U\)，ug 有定义；
\item
  \(vg^{-1}\) 对所有 \(v \in V\) 都有定义；
\item
  映射 \(u \mapsto ug\)、\(v \mapsto vg^{-1}\) 彼此互为从 U 到 V 及从 V 到 U 的解析同构。
\end{enumerate}

  由于乘积的定义域是 \(\mathbf{G} \times \mathbf{G}\) 中的开集，存在 \(e\) 的开邻域 U 和 \(g\) 的开邻域 V，具有性质 (a) 和 (b)。令 \(\eta(u) = ug\)（\(u \in \mathrm{U}\)），\(\eta'(v) = vg^{-1}\)（\(v \in \mathrm{V}\)）。缩小 U 和 V，可设 \((ug)g^{-1} = u\) 以及 \((vg^{-1})g = v\)，其中 \(u \in \mathrm{U}\) 和 \(v \in \mathrm{V}\)。于是 \(\eta\) 和 \(\eta'\) 是单射。再缩小 U，可设 \(\eta(\mathrm{U}) \subset \mathrm{V}\)。于是 \(\eta'(\mathrm{V}) \supset \mathrm{U}\)，且 \(\eta(\mathrm{U})\) 是 \(\eta'\) 下 U 的逆像，因而是 V 中 \(g\) 的开邻域。将 V 替换为 \(\eta(\mathrm{U})\)，最终得到 \(\eta\) 和 \(\eta'\) 是解析的互逆双射。

设 \(G_{1}, G_{2}\) 是两个李群芽，单位元分别为 \(e_{1}, e_{2}\)。从 \(G_{1}\) 到 \(G_{2}\) 的映射 f 若满足下列条件，就称为态射：

\begin{enumerate}
\def\labelenumi{(\roman{enumi})}
\item
  \(f\) 是解析的；
\item
  \(f(e_{1}) = e_{2};\)
\item
  若 \(g, h\) 是 \(\mathbf{G}_1\) 中满足 \(gh\) 有定义的元素，则 \(f(g)f(h)\) 有定义且等于 \(f(gh)\)。
\end{enumerate}

令 \(g \in G_1\)。由于 \(gg^{-1}\) 有定义且等于 \(e_1, f(g)f(g^{-1})\) 有定义并等于 \(e_2\)，从而

\[
f (g) ^ {- 1} = f (g) ^ {- 1} \left(f (g) f \left(g ^ {- 1}\right)\right) = \left(f (g) ^ {- 1} f (g)\right) f \left(g ^ {- 1}\right)
\]

即

\[
f (g) ^ {- 1} = f (g ^ {- 1}).\tag{11}
\]

两个态射的复合仍是态射。

若 \(f: \mathbf{G}_1 \to \mathbf{G}_2\) 和 \(f': \mathbf{G}_2 \to \mathbf{G}_1\) 互为逆态射，则它们是同构（特别使用公式 (11)）。

设 \(G_{1}\)、\(G_{2}\) 是两个李群芽，f 是从 \(G_{1}\) 到 \(G_{2}\) 的映射，满足上述条件 (ii) 和 (iii)，并且在 \(e_{1}\) 的一个开邻域内解析。利用命题 20，可仿照命题 4 证明 f 是态射。

设 \((\mathbf{G}, e, \theta, m)\) 是李群芽，\(\Omega\) 是 \(m\) 的定义集。令 \(\mathbf{H}\) 为 \(\mathbf{G}\) 的子流形，它包含 \(e\) 且在 \(\theta\) 下稳定。假设集合 \(\Omega_1\)（由 \((x, y) \in \Omega \cap (\mathbf{H} \times \mathbf{H})\) 中满足 \(m(x, y) \in \mathbf{H}\) 的元素所成）在 \(\mathbf{H} \times \mathbf{H}\) 中是开集。则 \((\mathbf{H}, e, \theta | \mathbf{H}, m | \Omega_1)\) 是李群芽。这样的李群芽称为 \(\mathbf{G}\) 的李子群芽。将 \(\mathbf{H}\) 典范嵌入 \(\mathbf{G}\) 的映射是态射。若 \(f: \mathbf{L} \to \mathbf{G}\) 是李群芽的态射且 \(f(\mathbf{L}) \subset \mathbf{H}\)，则 \(f: \mathbf{L} \to \mathbf{H}\) 是李群芽的态射。

设 K 的特征为 0。若将 H 是 G 的子流形这一假设替换为 H 是 G 的拟子流形这一假设，上述段落的结果仍然成立（参见《Differentiable and Analytic Manifolds》R, 5.8.5）。此时 H 称为 G 的李拟子群芽。

若 G 是单位元为 e 的李群芽，则 G 中 e 的每个对称开邻域都是 G 的李子群芽。（当 G 是李群时尤其如此。）设 H 是 G 的李子群芽；若 H 是 G 中 e 的邻域，则根据命题 20，H 在 G 中是开集。

有限个李群芽的乘积李群芽按显然的方式定义。

命题 21. 设 G、H 是两个李群芽，\(\phi\) 是从 G 到 H 的态射。下列条件等价：

\begin{enumerate}
\def\labelenumi{(\roman{enumi})}
\item
  \(\phi\) 在 \(e\) 处是 étale 的；
\item
  存在 G、H 的开李子群芽 G'、H'，使得 \(\phi|G'\) 是从 G' 到 H' 的同构。
\end{enumerate}

蕴含 \(\langle \mathrm{ii}\rangle \Rightarrow \langle \mathrm{i}\rangle\) 是显然的。假设 \(\phi\) 在 \(e\) 处是 étale 的。存在开李子群芽 \(\mathbf{G}_1\)，它是 \(\mathbf{G}\) 的开李子群芽，使得 \(\phi (\mathbf{G}_1)\) 在 \(\mathbf{H}\) 中是开集，并且 \(\phi |\mathbf{G}_1\) 是流形 \(\mathbf{G}_1\) 到流形 \(\phi (\mathbf{G}_1)\) 的同构。于是存在开李子群芽 \(\mathbf{G}'\)，它是 \(\mathbf{G}_1\) 的开李子群芽，使得 \(\mathbf{G}\) 中 \(\mathbf{G}'\) 的两个元素的乘积总有定义且属于 \(\mathbf{G}_1\)。若 \(g, g'\) 是 \(\mathbf{G}'\) 中满足 \(\mathrm{gg}^{\prime}\in \mathrm{G}'\) 的元素，则 \(\phi (g)\phi (g') = \phi (gg')\in \phi (\mathbf{G}')\)；若 \(g, g'\) 是 \(\mathbf{G}'\) 中满足 \(gg^{\prime}\in G_1 - G'\) 的元素，则

\[
\phi (g) \phi (g ^ {\prime}) = \phi (g g ^ {\prime}) \in \phi (\mathrm{G} _ {1}) - \phi (\mathrm{G} ^ {\prime}).
\]

因此 \(\phi |G^{\prime}\) 是李群芽 \(G^{\prime}\) 到 H 的开李子群芽 \(\phi (G^{\prime})\) 的同构。

若命题 21 的条件成立，则称 G 和 H 局部同构。

命题 22. 设 H 是李群，U 是 H 的李子群芽，N 是由 \(g \in \mathrm{H}\) 所成的集合，其中 U 与 \(g\mathrm{Ug}^{-1}\) 在 \(e\) 处具有相同的芽（《General Topology》，第 I 章，第 § 6 节，编号 10）。则 \(N\) 是包含 U 的 H 的子群。在 \(\mathbb{N}\) 上存在唯一一个具有下列性质的解析流形结构：

\begin{enumerate}
\def\labelenumi{(\roman{enumi})}
\item
  带此结构的 N 是李群；
\item
  U 是 N 的开子流形；
\item
  从 \(\mathbf{N}\) 到 \(\mathbf{H}\) 的典范嵌入是浸入。
\end{enumerate}

显然，\(\mathbf{N}\) 是 \(\mathbf{H}\) 的子群。若 \(g \in \mathbf{U}\)，则 \(ge \in \mathbf{U}\) 且 \(geg^{-1} \in \mathbf{U}\)，因而 \(gu \in \mathbf{U}\) 且 \(gug^{-1} \in \mathbf{U}\)，其中 \(u\) 充分接近 \(e\)（在 \(\mathbf{U}\) 中），所以 \(g\mathrm{U}g^{-1}\) 在 \(e\) 处的芽包含于 \(\mathbf{U}\) 的芽；交换 \(g\) 和 \(g^{-1}\) 后，可见 \(g\mathrm{U}g^{-1}\) 和 \(\mathbf{U}\) 在 \(e\) 处的芽相等。因此 \(\mathbf{U} \subset \mathbf{N}\)。

令 \(\mathbf{V}\) 为 \(e\) 在 \(\mathbf{U}\) 中的开邻域，使 \(\mathrm{V} = \mathrm{V}^{-1}\)、\(\mathrm{V}^2 \subset \mathrm{U}\)。第 9 号命题 18 的条件 (i)、(ii)、(iii)（其中以 \(\mathrm{G}\) 替换为 \(\mathrm{N}\)）成立。于是，在 \(\mathbf{N}\) 上存在解析流形结构，具有下列性质：\((\alpha)\) 带此结构的 \(\mathbf{N}\) 是李群；\((\beta)\) \(\mathbf{V}\) 在 \(\mathbf{N}\) 中是开集；\((\gamma)\) \(\mathbf{N}\) 和 \(\mathbf{U}\) 上的流形结构在 \(\mathbf{V}\) 上诱导出相同的结构。由于 \(\mathbf{V}\) 是 \(\mathbf{H}\) 的子流形，从 \(\mathbf{N}\) 到 \(\mathbf{H}\) 的典范嵌入在 \(e\) 处是浸入，因而在 \(\mathbf{N}\) 的每一点处都是浸入。令 \(u \in \mathbf{U}\)。存在 \(V'\)，它是 e 在 V 中的开邻域，使映射 \(v \mapsto vu\) 在 \(V'\) 上是解析同构，映到 U 中 u 的一个开邻域（命题 20），同时也映到 N 中 u 的一个开邻域。因此 U 在 N 中是开集，并且 U 的恒等映射关于 U 上给定的流形结构和 N 上的开子流形结构是同构；换言之，U 是 N 的开子流形。

最后，考虑 N 上具有命题条件 (i)、(ii) 的解析流形结构，并令 \(N^{*}\) 为由此得到的李群。于是从 N 到 \(N^{*}\) 的恒等映射在 e 处是 étale 的，因而是李群同构。这证明了命题的唯一性断言。

设 H 是李群，U 是 H 的李拟子群芽，N 是由 \(g \in H\) 所成的集合，其中 U 与 \(gUg^{-1}\) 在 e 处具有相同芽。若 K 的特征为 0，则在 G 上存在唯一一个具有命题 22 的条件 (i)、(ii) 的流形结构。证明与命题 22 相同。

推论。沿用命题 22 的记号，令 G 为 H 中由 U 生成的子群。则 G 是 N 的开子群。在 G 上存在唯一一个李群结构，使 U 是 G 的开子流形，且从 G 到 H 的典范嵌入是浸入。

注。沿用命题 22 及其推论的记号，假设 K 的特征为 0，H 是有限维的，且 U 上的拓扑具有可数基。即使有这些假设，也可能 G 在 H 中不是闭集（习题 3）。但是，若 G 是闭集，则 G 是 H 的李子群。因为映射 \((g,h)\mapsto gh\) 是 G 在 H 上的解析左作用律，e 的轨道是 G。于是我们的断言由命题 2 (iv) 和命题 14 (iii) 得出。

